\begin{solapartat}
\punts{0,5} El flux del camp magnètic és el producte escalar del camp magnètic per l'àrea del circuit. Calculem primer l'àrea del circuit en funció del temps: $A(t) = 0,25(0,3 - 5,0\,t)$.

El camp magnètic és perpendicular al circuit; per tant, paral·lel al vector de superfície i, en conseqüència, $\phi(t) = B\cdot A(t) = 2\cdot 0,25(0,3 - 5,0\,t) = (0,15 - 2,5\,t)\mathrm{Wb}$.

\punts{0,35} La força electromotriu és $\varepsilon = -\frac{d\phi}{dt} = 2,5\ \mathrm{V}$.

\punts{0,4} El corrent induït sempre és contrari a la variació de flux de camp magnètic. En aquest cas, es redueix el flux i, per tant, el corrent induït crearà un camp magnètic en la mateixa direcció que el camp ja existent. Conseqüentment, la intensitat induïda gira en sentit antihorari.

\figura[0.35]{sol_fig1.png}
\end{solapartat}

\begin{solapartat}
\punts{0,25} Aplicant la llei d'Ohm, trobem la intensitat que circula pel circuit deguda a la força electromotriu induïda: $I = \frac{\varepsilon}{R} = 2,5/50 = 0,05\ \mathrm{A} = 50\ \mathrm{mA}$.

\punts{0,5} La força magnètica que actua sobre la barra és $\vec{F} = I\vec{l}\times\vec{B}$. La longitud de la barra i el camp magnètic són perpendiculars; per tant, el mòdul de la força magnètica serà: $F = IlB = 0,05\cdot 0,25\cdot 2 = 0,025\ \mathrm{N}$.

\punts{0,5} La direcció de la força és perpendicular a la barra i al camp i la regla de la mà dreta ens indica que és cap amunt en el pla del paper.

També es pot justificar a partir de la llei de Lenz: la força ha d'actuar de tal manera que s'oposi al moviment; com que es mou cap avall, la força apunta cap amunt.

\figura[0.4]{sol_fig2.png}
\end{solapartat}
